\BourbakiExerciseGroup

\begin{enumerate}
\def\labelenumi{\arabic{enumi})}
\item
  有限群 G 的与群结构相容的每个拓扑，可以这样得到：取包含 G 的一个正规子群 H 的那些集合作为单位元的邻域。
\item
半拓扑群是指一个群 \(\mathbf{G}\) ，带有一个拓扑，使得对每个 \(a \in \mathbf{G}\) ，平移 \(x \to ax\) 与 \(x \to xa\) 都在 \(\mathbf{G}\) 上连续，且对称 \(x \to x^{-1}\) 在 \(\mathbf{G}(*)\) 上连续。
\end{enumerate}

\begin{enumerate}
\def\labelenumi{\alph{enumi})}
\tightlist
\item
  一个滤子 \(\mathfrak{B}\) （在群 \(\mathbf{G}\) 上）是 \(e\) 的邻域滤子（在使 \(\mathbf{G}\) 成为半拓扑群的一个拓扑中），当且仅当 \(\mathfrak{B}\) 满足公理 \((\mathrm{GV}_{\mathrm{I}})\) 与 \((\mathrm{GV}_{\mathrm{III}})\) ，且滤子族
\end{enumerate}

\[
\mathfrak {B} (x) = x. \mathfrak {B} = \mathfrak {B}. x \quad (x \in G)
\]

满足第一章 § 1 第 2 节的公理 \((\mathbf{V}_{\mathbf{IV}})\) 。

\begin{enumerate}
\def\labelenumi{\alph{enumi})}
\setcounter{enumi}{1}
\tightlist
\item
  若 G 是无限群，试证：以 \(\varnothing\) 与 G 的有限子集的余集为开集的那个拓扑，使 G 成为一个非豪斯多夫的半拓扑群，其中 \(\{e\}\) 是闭的，但与 G 的群结构不相容。
\end{enumerate}

\(^{*}c)\) 在 R 上定义一个比通常拓扑更细的拓扑，使 R 成为一个半拓扑群但不是拓扑群{[}考虑一个递减数列 \((r_{n})\) （诸数 \textgreater0 ，趋于 0），并取诸对称区间 {]}-a, a{[} 与点集 \(\pm r_{n}\) 的余集的交）。 *

(*) 若 G 是局部紧的，这些条件蕴含 G 的拓扑与其群结构相容（第十章，§ 3，习题 25）。

\begin{enumerate}
\def\labelenumi{\arabic{enumi})}
\setcounter{enumi}{2}
\item
  设 A 是半拓扑（相应地，拓扑）群 G 的一个子集。若 V 取遍 G 中 e 的邻域滤子，则诸集合 AV 的交、或诸集合 VA （或诸集合 VAV）的交是 A 的闭包 \(\overline{A}\) 。
\item
  仿拓扑群是指一个群 G ，带有一个满足公理 \((\mathrm{GT}_{\mathbf{I}})(*)\) 的拓扑。
\end{enumerate}

\begin{enumerate}
\def\labelenumi{\alph{enumi})}
\item
  一个滤子 \(\mathfrak{B}\) （群 \(G\) 上的）是 \(e\) 的邻域滤子（对应于使 \(G\) 成为仿拓扑群的一个拓扑），当且仅当 \(\mathfrak{B}\) 满足公理 \((\mathrm{GV}_{\mathbf{I}})\) 与 \((\mathrm{GV}_{\mathbf{III}})\) ，且 \(e \in V\) 对一切 \(V \in \mathfrak{B}\) 成立。相应的拓扑则是唯一的；它是豪斯多夫的，当且仅当诸集合 \(V.V^{-1}\) （\(V\) 取遍 \(\mathfrak{B}\) ）的交仅由点 \(e\) 组成。
\item
  在整数群 \(\mathbf{Z}\) 上，设 \(\mathbf{V}_n\) 是由 \(0\) 与所有整数 \(m \geqslant n\) 组成的集合（对每个 \(n > 0\) ）。试证诸 \(\mathbf{V}_n\) 构成 \(0\) 的邻域基本系（对应 \(\mathbf{Z}\) 上一个非豪斯多夫拓扑），且带此拓扑的 \(\mathbf{Z}\) 是一个仿拓扑群，其中 \(\{0\}\) 是闭的。
\item
  仿拓扑群 G 是拓扑群，当且仅当对 G 的每个子集 A ，诸集合 AV （或诸集合 VA）的交是 A 的闭包 \(\overline{A}\) （当 V 取遍 e 的邻域滤子时）。（为看出条件是充分的，取 A 为 e 的一个开邻域的余集。）
\end{enumerate}

\begin{enumerate}
\def\labelenumi{\arabic{enumi})}
\setcounter{enumi}{4}
\tightlist
\item
  设 \(\mathfrak{B}\) 是群 \(\mathbf{G}\) 上的一个滤子，满足公理 \((\mathrm{GV}_{\mathbf{I}})\) 与 \((\mathrm{GV}_{\mathbf{II}})\) 。
\end{enumerate}

\begin{enumerate}
\def\labelenumi{\alph{enumi})}
\tightlist
\item
  存在唯一的拓扑 \(\mathfrak{C}_s\) （相应地，\(\mathfrak{C}_d\) ），它定义在 \(\mathbf{G}\) 上，使得对每个 \(a \in \mathbf{G}\) ，\(a\) 关于 \(\mathfrak{C}_s\) （相应地，\(\mathfrak{C}_d\) ）的邻域滤子是 \(a\) . \(\mathfrak{B}\) （相应地，\(\mathfrak{B}\) . \(a\) ）。
\end{enumerate}

设 \(\mathbf{G}_s\) （相应地，\(\mathbf{G}_d\) ）是通过给 \(\mathbf{G}\) 赋以拓扑 \(\mathcal{C}_s\) （相应地，\(\mathcal{C}_d\) ）所得到的拓扑空间。对每个 \(a \in \mathbf{G}\) ，左平移 \(x \to ax\) （相应地，右平移 \(x \to xa\) ）是 \(\mathbf{G}_s\) 到 \(\mathbf{G}_s\) 中（相应地，\(\mathbf{G}_d\) 到 \(\mathbf{G}_d\) 中）的连续映射。对称 \(x \to x^{-1}\) 是 \(\mathbf{G}_s\) 到 \(\mathbf{G}_d\) 中以及 \(\mathbf{G}_d\) 到 \(\mathbf{G}_s\) 中的连续映射。

\begin{enumerate}
\def\labelenumi{\alph{enumi})}
\setcounter{enumi}{1}
\tightlist
\item
  下面的条件等价：1) V 满足 \((\mathrm{GV}_{\mathrm{III}})\) ；
\end{enumerate}

\begin{enumerate}
\def\labelenumi{\arabic{enumi})}
\setcounter{enumi}{1}
\tightlist
\item
  对每个 \(a \in G, x \to xa\) 是 \(G_{s}\) 到 \(G_{s}\) 中的连续映射；
\item
  \(x \to x^{-1}\) 是 \(G_{s}\) 到 \(G_{s}\) 中的连续映射。
\end{enumerate}

\(^{*}c)\) 设 \(G\) 是群 \(\mathbf{GL}(2, \mathbf{Q}_p)\) （即 \(2 \times 2\) 方阵在 \(p\) 进域 \(\mathbf{Q}_p\) 上所成的群）。对每个整数 \(n > 0\) ，设 \(H_n\) 是形如 \(I + p^n U\) 的矩阵组成的集合，其中 \(U\) 是一个 \(2 \times 2\) 矩阵，其元素是 \(p\) 进整数。试证诸 \(H_n\) 是 \(G\) 的子群，且它们的交仅由单位元组成。推断：若 \(\mathfrak{B}\) 是以诸 \(H_n\) 为基的滤子，则相应的拓扑 \(T_s\) 与 \(T_d\) （在 \(G\) 上）是豪斯多夫的。试证这些拓扑是不同的。*

\begin{enumerate}
\def\labelenumi{\arabic{enumi})}
\setcounter{enumi}{5}
\tightlist
\item
  设 G 是拓扑群，V 是 e 在 G 中的一个开邻域。试证：\(x \in G\) （满足 \(x^{2} \in V\) 者）组成的集合，是使得 \(W^{2} \subset V\) 的 e 的邻域 W 的并。
\end{enumerate}

*7) 设 \(\varphi\) 是 \(\mathbf{R}\) 到 \(\mathbf{T} = \mathbf{R} / \mathbf{Z}\) 上的典范映射，并设 \(f\) 是映射 \(x \to \varphi(3x)\) 限制于邻域 \(\mathbf{V} = [-1/8, +1/8]\) （\(0\) 在 \(\mathbf{R}\) 中的）上所得的映射。试证 \(f\) 是 \(\mathbf{V}\) 到 \(f(\mathbf{V})\) 上的一个同胚，且满足第 3 节定义 2 的条件 1)，但不是 \(\mathbf{R}\) 到 \(\mathbf{T}\) 的局部同构。*

\begin{enumerate}
\def\labelenumi{\arabic{enumi})}
\setcounter{enumi}{7}
\item
设 G 是一个豪斯多夫拓扑群，并设 \(s \neq e\) 是 G 的一点。试证：存在 \(e\) 的一个对称邻域 V ，使得 \(s \notin V^2\) ，且 G 可以被至多 17 个形如 \(a_k V b_k (1 \leqslant k \leqslant 17)\) 的集合所覆盖。（在 \(e\) 的使得 \(s \notin V^2\) 的对称邻域 V 组成的集合中考虑一个极大元。这样我们就被引导去考虑 \(x \in G\) （满足 \(x^2 = s\) 者）组成的集合 S ；注意，若 \(x, y\) 是 S 的两个元素且 \(xy^{-1} \in S\) ，则必有 \(x^2 = e\) ，又若 \(z\) 是 S 的第三个元素，使得 \(xz^{-1} \in S\) 且 \(yz^{-1} \in S\) ，则 \(z\) 必与 \(xy^{-1}\) 交换，且 \(xy^{-1}z^{-1} \notin S\) 。）若 G 是交换的，数 17 可换成 5。
\item
  试证：在群 G 上，与 G 的群结构相容的一族拓扑的上确界与 G 的群结构相容。
\end{enumerate}

